\hypertarget{system-cost-optimization-migrating-from-apache-kafka-to-amazon-sqs}{%
\section{System Cost Optimization: Migrating from Apache Kafka to Amazon
SQS}}

A crucial criterion when designing systems on the Cloud is \textbf{Cost
Optimization}. In the first version (v1) developed in our local
environment, the project used \textbf{Apache Kafka} as the message
broker to stream data from the Crawler to the Database. However, when
migrating to AWS, cost became a significant barrier.

\hypertarget{the-cost-challenge-with-kafka}{%
\subsubsection{The Cost Challenge with
Kafka}}

Apache Kafka is an incredibly powerful stream processing tool. But to
run Kafka on AWS, we generally have two main choices: 1. \textbf{Amazon
MSK (Managed Streaming for Apache Kafka)}: Very expensive (can cost
dozens or hundreds of USD per month). 2. \textbf{Self-hosting Kafka on
EC2}: Requires maintaining EC2 instances running 24/7, incurring fixed
costs (at least \textasciitilde\$10-15/month) plus the overhead of
system administration, security updates, and managing ZooKeeper.

Meanwhile, the data flow of NewsRAG is fundamentally \textbf{Batch
processing}: The Crawler only runs 1-2 times a day (at night) and pushes
around 500 - 1000 new articles. Maintaining a Kafka cluster running 24/7
just to serve this intermittent data flow is extremely wasteful.

\hypertarget{the-alternative-amazon-sqs-simple-queue-service}{%
\subsubsection{The Alternative: Amazon SQS (Simple Queue
Service)}}

We decided to completely replace the queue architecture with
\textbf{Amazon SQS}.

\textbf{Why SQS is the perfect fit:} 1. \textbf{Fully Serverless \&
Pay-as-you-go}: SQS requires no servers. You only pay based on the
number of requests. With our system's daily news volume, the number of
requests easily fits within the \textbf{Free Tier} (the first 1 million
requests per month are free). The queue cost was reduced to
\textbf{essentially \$0}. 2. \textbf{Native Integration with AWS
Lambda}: SQS can act as an ``Event Source'' to automatically trigger a
Lambda function (Lambda Consumer). When the Crawler pushes an article to
SQS, SQS automatically invokes Lambda to save it into the Database,
without us having to write polling loops. 3. \textbf{Dead-Letter Queue
(DLQ)}: SQS provides a built-in DLQ mechanism. If an article is
malformed and Lambda fails to save it to the DB after several retries,
that message is moved to the DLQ so we can analyze it later, ensuring no
data loss.

\hypertarget{design-evaluation}{%
\subsubsection{Design Evaluation}}

By removing Kafka and switching to SQS, the system became not only
lighter and easier to maintain but also saved a massive amount in
operational budget. This is a practical proof that: \textbf{Choosing
technology is not about picking the ``fanciest'' tool, but choosing the
most ``appropriate'' one for the current scale and problem.}
